\documentclass{article}
\usepackage{iftex}
\ifXeTeX
  \usepackage{fontspec}
  \usepackage{xeCJK}
  \usepackage[haranoaji]{zxjafont}
  \XeTeXlinebreaklocale ""
\fi
\ifLuaTeX
  \usepackage[haranoaji]{luatexja-preset}
\fi
\usepackage[a4paper, margin=2cm]{geometry}
\usepackage{graphicx}
\newdimen\charwd
{\tt\global\setbox0=\hbox{x}\global\charwd=\wd0}
\frenchspacing
\usepackage{fancyhdr}
\pagestyle{fancy}
\renewcommand{\headrulewidth}{0pt}
\fancyhf{}
\fancyhead[R]{\rm src2tex-go version 2.23}
% plain TeX compatibility
\makeatletter
\providecommand{\eqalign}[1]{\vcenter{\openup1\jot\m@th
  \ialign{\strut\hfil$\displaystyle{##}$&$\displaystyle{{}##}$\hfil\crcr#1\crcr}}}
\makeatother
\fancyfoot[R]{\rm\hfill samples{\tt /}simpson.c\qquad page \thepage}
\begin{document}

\ifx\sevenrm\undefined
  \font\sevenrm=cmr7 scaled \magstep0
\fi

\newread\MyStyle
\openin\MyStyle=src2latex.s2t
\ifeof\MyStyle
  \closein\MyStyle
\else
  \input src2latex.s2t
  \closein\MyStyle
\fi

\def\mc{\relax}
\def\gt{\relax}
\def\sc{\scshape}

\tt\mc 

\noindent
\rm\mc {\tt /}{\tt *}\kern1\charwd  \hrulefill \rm\mc \kern1\charwd {\tt *}

\noindent
\mbox{}\hfill

\noindent
\mbox{} \bf Simpson \rm\mc の公式

\noindent
\mbox{}\hfill

\noindent
\mbox{} {}% beginning of TeX mode

\noindent
定積分 $\displaystyle\int_a^b f(x)\,dx$ を計算する一つの方法として、
次の近似式が良く使われる。

$$\int_a^b f(x)\,dx\sim{h\over3}\,(y_0+4y_1+2y_2+4y_3+2y_4+4y_5+
					\cdots +2y_{n-2}+4y_{n-1}+y_n)\ ,$$

\noindent
ここで、自然数 $n$ は偶数とし、
$\displaystyle x_i=a+{i\over n}\,(b-a),\ y_i=f(x_i)\quad (i=0, 1, 2, \cdots, n)$
とする。
この公式の証明は、比較的やさしい。
じっさい、Taylor の定理から導かれる等式

$$\int_{\xi-{1\over h}}^{\xi+{1\over h}}f(x)\,dx
	\sim\int_{\xi-{1\over h}}^{\xi+{1\over h}}p(x)\,dx
	={h\over3}\,\bigl(f(\xi-{1\over h})+4f(\xi)+f(\xi+{1\over h})\bigr)+O(h^5)$$

\noindent
を $\,\displaystyle{n\over2}\,$ 個足し合わせることにより、
この公式は証明される。
\begin{center}
\includegraphics[scale=.7]{simpson}
\end{center}


\noindent
ここで、$p(x)$ は次のような条件を満足する２次の多項式を表す:
$$
p\Bigl(\xi-{1\over h}\Bigr)=f\Bigl(\xi-{1\over h}\Bigr),
\ p(\xi)=f(\xi),
\ p\Bigl(\xi+{1\over h}\Bigr)=f\Bigl(\xi+{1\over h}\Bigr)\ .
$$

\noindent
{\tt A, B, F(X), N} の定義をいろいろと変えて、各人で数値実験を行ってみよう。

% end of TeX mode  \rm\mc 

\noindent
\mbox{}\hfill

\noindent
\mbox{}{\tt *}\kern1\charwd  \hrulefill \rm\mc \kern1\charwd {\tt *}{\tt /}
\tt\mc 

\noindent
\mbox{}\hfill

\noindent
\mbox{}\rm\mc {\tt /}{\tt *}\kern1\charwd  \hrulefill\ simpson.c\ \hrulefill \rm\mc \kern1\charwd {\tt *}{\tt /}
\tt\mc 

\noindent
\mbox{}\hfill

\noindent
\mbox{}\hfill

\noindent
\mbox{}{\tt\#}include{\tt\mc \kern1\charwd}{\tt <}stdio.h{\tt >}

\noindent
\mbox{}{\tt\#}define{\tt\mc \kern1\charwd}A{\tt\mc \kern1\charwd}0.{\tt\mc \kern4\charwd}{\tt\mc \kern8\charwd}{\tt\mc \kern8\charwd}{\tt\mc \kern8\charwd}\rm\mc {\tt /}{\tt *}\kern1\charwd 積分をする区間\kern1\charwd $[a, b]$\rm\mc \kern1\charwd  \hfill \rm\mc \kern1\charwd {\tt *}{\tt /}
\tt\mc 

\noindent
\mbox{}{\tt\#}define{\tt\mc \kern1\charwd}B{\tt\mc \kern1\charwd}1.

\noindent
\mbox{}{\tt\#}define{\tt\mc \kern1\charwd}F(X){\tt\mc \kern1\charwd}((X){\tt *}(X)){\tt\mc \kern2\charwd}{\tt\mc \kern8\charwd}{\tt\mc \kern8\charwd}\rm\mc {\tt /}{\tt *}\kern1\charwd 被積分関数\kern1\charwd $f(x)$\rm\mc \kern1\charwd  \hfill \rm\mc \kern1\charwd {\tt *}{\tt /}
\tt\mc 

\noindent
\mbox{}{\tt\#}define{\tt\mc \kern1\charwd}N{\tt\mc \kern1\charwd}20{\tt\mc \kern4\charwd}{\tt\mc \kern8\charwd}{\tt\mc \kern8\charwd}{\tt\mc \kern8\charwd}\rm\mc {\tt /}{\tt *}\kern1\charwd  {}区間 $[a, b]$ の分割数 $n$
					   \hfill \rm\mc \kern1\charwd {\tt *}{\tt /}
\tt\mc 

\noindent
\mbox{}\hfill

\noindent
\mbox{}double{\tt\mc \kern1\charwd}simpson{\tt\_\kern.141em}rule(){\tt\mc \kern3\charwd}{\tt\mc \kern8\charwd}{\tt\mc \kern8\charwd}\rm\mc {\tt /}{\tt *}\kern1\charwd  {}Simpson の公式を用いて
					   $\displaystyle\int_a^b f(x)\,dx$
					   を計算する関数 \hfill \rm\mc \kern1\charwd {\tt *}{\tt /}
\tt\mc 

\noindent
\mbox{}{\tt\char'173}

\noindent
\mbox{}{\tt\mc \kern4\charwd}long{\tt\mc \kern1\charwd}i;{\tt\mc \kern5\charwd}{\tt\mc \kern8\charwd}{\tt\mc \kern8\charwd}{\tt\mc \kern8\charwd}\rm\mc {\tt /}{\tt *}\kern1\charwd long\kern1\charwd 整数を使う\kern1\charwd  \hfill \rm\mc \kern1\charwd {\tt *}{\tt /}
\tt\mc 

\noindent
\mbox{}{\tt\mc \kern4\charwd}double{\tt\mc \kern1\charwd}h,{\tt\mc \kern1\charwd}s,{\tt\mc \kern1\charwd}y[N{\tt\mc \kern1\charwd}+{\tt\mc \kern1\charwd}1];{\tt\mc \kern6\charwd}{\tt\mc \kern8\charwd}\rm\mc {\tt /}{\tt *}\kern1\charwd  {}メッシュ $h$, 積分値 $s$,
					   分点 $\displaystyle{i\over n}$
					   での関数値 $y_i$ \hfill \rm\mc \kern1\charwd {\tt *}{\tt /}
\tt\mc 

\noindent
\mbox{}\hfill

\noindent
\mbox{}{\tt\mc \kern4\charwd}h{\tt\mc \kern1\charwd}={\tt\mc \kern1\charwd}1.{\tt\mc \kern1\charwd}{\tt /}{\tt\mc \kern1\charwd}(double){\tt\mc \kern1\charwd}N;{\tt\mc \kern8\charwd}{\tt\mc \kern8\charwd}\rm\mc {\tt /}{\tt *}\kern1\charwd  {}$\displaystyle h={1\over n}$
					 \hfill \rm\mc \kern1\charwd {\tt *}{\tt /}
\tt\mc 

\noindent
\mbox{}{\tt\mc \kern4\charwd}for{\tt\mc \kern1\charwd}(i{\tt\mc \kern1\charwd}={\tt\mc \kern1\charwd}0;{\tt\mc \kern1\charwd}i{\tt\mc \kern1\charwd}{\tt <}={\tt\mc \kern1\charwd}N;{\tt\mc \kern1\charwd}++i){\tt\mc \kern4\charwd}{\tt\mc \kern8\charwd}\rm\mc {\tt /}{\tt *}\kern1\charwd  {}$\displaystyle y_i
					   =f\bigl({i\over n}\bigr)
					   \quad (i=0, 1, 2, \cdots, n)$
					   \hfill \rm\mc \kern1\charwd {\tt *}{\tt /}
\tt\mc 

\noindent
\mbox{}{\tt\mc \kern8\charwd}y[i]{\tt\mc \kern1\charwd}={\tt\mc \kern1\charwd}F((double){\tt\mc \kern1\charwd}i{\tt\mc \kern1\charwd}{\tt /}{\tt\mc \kern1\charwd}(double){\tt\mc \kern1\charwd}N);

\noindent
\mbox{}{\tt\mc \kern8\charwd}{\tt\mc \kern8\charwd}{\tt\mc \kern8\charwd}{\tt\mc \kern8\charwd}{\tt\mc \kern8\charwd}\rm\mc {\tt /}{\tt *}\kern1\charwd 以下は\kern1\charwd  \hfill \rm\mc \kern1\charwd {\tt *}{\tt /}
\tt\mc 

\noindent
\mbox{}{\tt\mc \kern8\charwd}{\tt\mc \kern8\charwd}{\tt\mc \kern8\charwd}{\tt\mc \kern8\charwd}{\tt\mc \kern8\charwd}\rm\mc {\tt /}{\tt *}\kern1\charwd $\displaystyle
					   s={h\over3}(y_0+4y_1+2y_2+4y_3+
					   \cdots +4y_{n-1}+y_n)$\rm\mc \kern1\charwd  \hfill \rm\mc \kern1\charwd {\tt *}{\tt /}
\tt\mc 

\noindent
\mbox{}{\tt\mc \kern8\charwd}{\tt\mc \kern8\charwd}{\tt\mc \kern8\charwd}{\tt\mc \kern8\charwd}{\tt\mc \kern8\charwd}\rm\mc {\tt /}{\tt *}\kern1\charwd の計算\kern1\charwd  \hfill \rm\mc \kern1\charwd {\tt *}{\tt /}
\tt\mc 

\noindent
\mbox{}{\tt\mc \kern4\charwd}s{\tt\mc \kern1\charwd}={\tt\mc \kern1\charwd}y[0];

\noindent
\mbox{}{\tt\mc \kern4\charwd}for{\tt\mc \kern1\charwd}(i{\tt\mc \kern1\charwd}={\tt\mc \kern1\charwd}1;{\tt\mc \kern1\charwd}i{\tt\mc \kern1\charwd}{\tt <}{\tt\mc \kern1\charwd}N;{\tt\mc \kern1\charwd}++i)

\noindent
\mbox{}{\tt\mc \kern4\charwd}{\tt\char'173}

\noindent
\mbox{}{\tt\mc \kern8\charwd}if{\tt\mc \kern1\charwd}(i{\tt\mc \kern1\charwd}{\tt\%}{\tt\mc \kern1\charwd}2{\tt\mc \kern1\charwd}=={\tt\mc \kern1\charwd}1){\tt\mc \kern1\charwd}{\tt\mc \kern8\charwd}{\tt\mc \kern8\charwd}\rm\mc {\tt /}{\tt *}\kern1\charwd  {}ただし {\tt i \% 2}= $i\ $
					   を 2 で割った余り \hfill \rm\mc \kern1\charwd {\tt *}{\tt /}
\tt\mc 

\noindent
\mbox{}{\tt\mc \kern8\charwd}{\tt\mc \kern4\charwd}s{\tt\mc \kern1\charwd}+={\tt\mc \kern1\charwd}4.{\tt\mc \kern1\charwd}{\tt *}{\tt\mc \kern1\charwd}y[i];

\noindent
\mbox{}{\tt\mc \kern8\charwd}else

\noindent
\mbox{}{\tt\mc \kern8\charwd}{\tt\mc \kern4\charwd}s{\tt\mc \kern1\charwd}+={\tt\mc \kern1\charwd}2.{\tt\mc \kern1\charwd}{\tt *}{\tt\mc \kern1\charwd}y[i];

\noindent
\mbox{}{\tt\mc \kern4\charwd}{\tt\char'175}

\noindent
\mbox{}{\tt\mc \kern4\charwd}s{\tt\mc \kern1\charwd}+={\tt\mc \kern1\charwd}y[N];

\noindent
\mbox{}{\tt\mc \kern4\charwd}s{\tt\mc \kern1\charwd}{\tt *}={\tt\mc \kern1\charwd}(h{\tt\mc \kern1\charwd}{\tt /}{\tt\mc \kern1\charwd}3.);

\noindent
\mbox{}{\tt\mc \kern4\charwd}return{\tt\mc \kern1\charwd}s;{\tt\mc \kern3\charwd}{\tt\mc \kern8\charwd}{\tt\mc \kern8\charwd}{\tt\mc \kern8\charwd}\rm\mc {\tt /}{\tt *}\kern1\charwd $s$\rm\mc \kern1\charwd の値を返す\kern1\charwd  \hfill \rm\mc \kern1\charwd {\tt *}{\tt /}
\tt\mc 

\noindent
\mbox{}{\tt\char'175}

\noindent
\mbox{}\hfill

\noindent
\mbox{}main()

\noindent
\mbox{}{\tt\char'173}

\noindent
\mbox{}{\tt\mc \kern4\charwd}double{\tt\mc \kern1\charwd}s;

\noindent
\mbox{}\hfill

\noindent
\mbox{}{\tt\mc \kern4\charwd}s{\tt\mc \kern1\charwd}={\tt\mc \kern1\charwd}simpson{\tt\_\kern.141em}rule();{\tt\mc \kern1\charwd}{\tt\mc \kern8\charwd}{\tt\mc \kern8\charwd}\rm\mc {\tt /}{\tt *}\kern1\charwd  {}上で定義した関数
					   simpson\_rule() を使う \hfill \rm\mc \kern1\charwd {\tt *}{\tt /}
\tt\mc 

\noindent
\mbox{}{\tt\mc \kern4\charwd}printf({\tt "}{\tt\%}.16f{\tt\char92}n{\tt "},{\tt\mc \kern1\charwd}s);{\tt\mc \kern7\charwd}{\tt\mc \kern8\charwd}\rm\mc {\tt /}{\tt *}\kern1\charwd  {}$s$ を小数点以下16桁表示する
					 \hfill \rm\mc \kern1\charwd {\tt *}{\tt /}
\tt\mc 

\noindent
\mbox{}{\tt\char'175}

\noindent
\mbox{}

\rm\mc

\end{document}
